%% Description of the HUST bearing data dataset.
%% Covers dimensions, combinations of faults/bearings/conditions, EDA findings,
%% supported tests and splits; the test rig is not described here.
%% Build (from eda/hust/):   pdflatex hust && pdflatex hust

\documentclass[11pt,a4paper]{article}
\usepackage[T1]{fontenc}
\usepackage[utf8]{inputenc}
\usepackage[margin=2.4cm]{geometry}
\usepackage{amsmath}
\usepackage{booktabs}
\usepackage{graphicx}
\usepackage[hidelinks]{hyperref}

\title{HUST bearing data}
\author{}
\date{\today}

\begin{document}
\maketitle

\section{Source}

\section{Dimensions of the time series}

\section{Combinations of faults, bearings and operating conditions}

\begin{table}[h]
\centering
\small
\begin{tabular}{l ccccc}
\toprule
Fault & 6204 & 6205 & 6206 & 6207 & 6208 \\
\midrule
N (healthy)        & 024 & 024 & 024 & 024 & 024 \\
I (inner)          & 024 & 024 & 024 & 024 & 024 \\
O (outer)          & 024 & 024 & 024 & 024 & 024 \\
B (ball)           & --  & 024 & 024 & 024 & 024 \\
IO (inner + outer) & 024 & 024 & 024 & 024 & 024 \\
IB (inner + ball)  & --  & 024 & 024 & 024 & 024 \\
OB (outer + ball)  & 024 & 024 & 024 & 024 & 024 \\
\bottomrule
\end{tabular}
\end{table}

\begin{itemize}
  \item Columns: bearing type (6204: 20\,mm bore \dots\ 6208: 40\,mm bore), so defect
        frequencies differ per column.
  \item Each digit is one recording at that load: 0, 200, 400\,W (powder brake; the load
        digit of the file name). 99 recordings; -- means no file (stated in the data note).
  \item One recording: 10\,s at 51.2\,kHz (512,000 samples), one vibration channel.
        Not all are 512,000 samples: 9 are shorter (shortest OB402, 121,082), some longer (D223).
  \item Faults are wire-cut cracks, 0.2\,mm wide; IO, IB, OB are two defects on one bearing.
  \item Whether each cell is one physical bearing is unclear: the data note speaks of 27
        defective and 3 healthy bearings, the files hold 28 fault/type combinations and 5
        healthy types (D223). The card marks the loads as inferred.
  \item Source: \texttt{FAULTS}, \texttt{BEARINGS}, \texttt{LOADS\_W} in \texttt{harness/data/hust.py}.
\end{itemize}

\section{Findings from EDA}

\section{Tests the data support}

\section{Splits}

\end{document}
